\documentclass[10pt]{article}
\usepackage[margin=20mm]{geometry}
\usepackage{amsmath,amssymb,mathtools,fontspec,unicode-math,xurl,hyperref}
\IfFontExistsTF{DejaVu Serif}{\setmainfont{DejaVu Serif}}{\setmainfont{Latin Modern Roman}}
\setmathfont{latinmodern-math.otf}
\setlength{\parindent}{0pt}\setlength{\parskip}{6pt}\sloppy
\title{Entropy normalization and variational consistency
in the porous medium equation}
\author{AI-assisted technical review}
\date{}
\begin{document}
\maketitle
Abstract: The acquired preprint derives a porous-medium hydrodynamic limit, a dynamical large-deviation principle, and an entropy–action identity from a rescaled zero-range process. The relative entropy printed in its definition lacks the reference-density weight used by the equilibrium law and subsequent calculations. Exact closed events contradict the resulting large-deviation upper bound, and a smooth positive path contradicts the stated entropy–action identity. A separate check of the central gradient-flow proof finds an infeasible proposed control and inconsistent action coefficients. Local discrepancies in a sharp discrete-to-continuum bound, a critical exponential moment, and a uniform mass estimate are also demonstrated. The highest supported assessment is Major. Explicit corrections are verified at the stated local scope; no complete revised irregular-path theory is certified. Scientific adjudication remains pending.\par
\section*{Identity, source and scope}
The reviewed paper is Benjamin Gess and Daniel Heydecker, The porous medium equation: Large deviations and gradient flow with degenerate and unbounded diffusion, acquired as arXiv:2303.11289v3, 24 March 2025, 62 pages [GH]. The associated publication is Communications on Pure and Applied Mathematics 78(9) (2025), 1609–1655, DOI `10.1002/cpa.22251`. References to source pages below mean the printed pages of this preprint; they agree with its PDF page indices counted from one.\par
The entire acquired source, including Appendices A–B and the bibliography, was read. Detailed checks concern the entropy normalization, equilibrium partition function, elementary large-deviation events, variational control algebra, discrete dissipation, and endpoint moment and mass estimates. The nonlinear PDE approximation and uniqueness results imported from the source’s reference [31], and all external tightness and discrete Sobolev theorems, were not independently reproved. Reading those dependencies is distinguished from certifying every theorem of the paper.\par
A bounded publisher-HTML comparison was made for version context; the journal PDF was not acquired and full version equivalence is not established. All counterexamples below are anchored to the retained v3 PDF, not to a claim of byte-identical journal text. No author correction or adoption of the corrections below is established by the inspected material.\par
\section*{Definitions used in the checks}
Equation reference: eq:printedH, eq:skeleton, eq:action, eq:continuousD. The torus has unit volume. Write \[h(s)=s\log s-s+1,\qquad h(0)=1.\] Source (3.6), p. 12, explicitly defines \[
\mathcal H_\rho(u)=\int_{\mathbb T^d}h(u/\rho)\,dx.\] The measure in that display is \(dx\), with no factor \(\rho(x)\). The definitions in source (1.6)–(1.8), p. 3, are \[
\partial_tu=\tfrac12\Delta u^\alpha-\nabla\cdot(u^{\alpha/2}g),
\qquad
\mathcal J(u)=\tfrac12\inf_g\|g\|_{L^2_{t,x}}^2,
\qquad
\mathcal I_\rho(u)=\alpha\mathcal H_\rho(u_0)+\mathcal J(u).\] The action in source (1.17)–(1.18), p. 5, is \[
\mathcal A(u)=\tfrac12\inf_\theta\|\theta\|_{L^2_{t,x}}^2,
\qquad
\partial_tu+\nabla\cdot(\tfrac12u^{\alpha/2}\theta)=0.\] The infima are over controls satisfying the displayed equations; these constraints cannot be changed when comparing their values. The dissipation in source (1.11), p. 4, is \[
\mathcal D_\alpha(u)=\frac2\alpha\int_{\mathbb T^d}|\nabla u^{\alpha/2}|^2\,dx.\]\par
\section*{Problem 1: the relative entropy gives the wrong rate}
Problem 1 (Major: reference-density weight).  With the literal definition [eq:printedH] (\#eq:printedH), the large-deviation upper bound in source Theorem 2 is false. The same normalization also invalidates the relative-entropy identity (1.19) of Theorem 4 for allowed constant reference densities.\par
\subsection*{An exact absorbing event}
Set \(d=1\), \(\alpha=2\), \(t_{\mathrm{fin}}=1\), \(\rho\equiv\tfrac12\), and \(\chi_N=N^{-2}\). All parameters are allowed: \(\rho\) is continuous and strictly positive, \(\alpha>1\), and the scaling condition (1.2), p. 2, is exactly \(N^2\chi_N=1\). Use the prescribed product initial law \(\Pi_\rho^N\).\par
Equation reference: eq:partition. Source §3.2, p. 13, gives \[
\pi_\rho^N(\chi_Nm)=
 \frac{(\rho/\chi_N)^{2m}}{(m!)^2 Z_2(\rho/\chi_N)},
\qquad
Z_2(\phi)=\sum_{m=0}^{\infty}\frac{\phi^{2m}}{(m!)^2}.\] The zero configuration is absorbing because every jump rate vanishes. Thus the singleton \(E_0=\{u_t\equiv0\}\), a closed subset of the stated Skorokhod space, has the exact probability \[\mathbb P(\eta^N\in E_0)=Z_2(N^2/2)^{-N}.\] The zero path belongs to the energy class, solves [eq:skeleton] (\#eq:skeleton) with \(g=0\), and has \(\mathcal J(0)=0\). Its printed rate is \(\mathcal I_{1/2}(0)=2\mathcal H_{1/2}(0)=2\).\par
Equation reference: eq:zeroContradiction. For completeness, the partition asymptotic needed here follows without any large-deviation theorem. Since its terms are nonnegative, \(Z_2(\phi)\le(\sum_m\phi^m/m!)^2=e^{2\phi}\). For \(m=\lfloor\phi\rfloor\) and \(\phi\ge2\), monotonicity of \(\log\) gives \[\log(m!)\le m\log m-m+1+\log m.\] The single \(m\)-term of the sum consequently gives \[2\phi-4-2\log\phi\le\log Z_2(\phi)\le2\phi.\] In particular, \(\phi^{-1}\log Z_2(\phi)\to2\), consistently with source Lemma 4.1, p. 17. Hence \[
\lim_{N\to\infty}\frac{\chi_N}{N}\log\mathbb P(\eta^N\in E_0)
=-\lim_{N\to\infty}\frac{\log Z_2(N^2/2)}{N^2}
=-1.\] The claimed closed-set upper bound (1.9) instead requires this limit to be at most \(-2\). This is a direct contradiction, not an inference from an unproved replacement lemma.\par
\subsection*{A strictly positive initial-profile witness}
The contradiction does not require admitting a zero-mass fluctuation. Let \(E_1=\{u\in\mathbb D:u_0\equiv1\}\), which is closed because evaluation at time zero is continuous. The exact configuration \(\eta_0^N\equiv1\) has \(m=N^2\) particles per site and positive total mass. Future dynamics are unrestricted in this event. Formula [eq:partition] (\#eq:partition) and the elementary factorial asymptotic \(\log(m!)=m\log m-m+o(m)\) yield \[\lim_N\frac{\chi_N}{N}\log\mathbb P(\eta^N\in E_1)=1-2\log2.\] On the other hand, \(\mathcal J\ge0\), and the constant macroscopic path \(u_t\equiv1\) attains \(\mathcal J=0\). Therefore \[\inf_{E_1}\mathcal I_{1/2}=2h(2)=4\log2-2.\] The printed upper bound would be \(2-4\log2\), strictly less than the actual value \(1-2\log2\), since \(\log2>\tfrac12\). This checks the same failure on a strictly positive, finite-entropy initial profile.\par
\subsection*{The mismatch in the dual and change-of-measure calculation}
Equation reference: eq:legendre. For \(u,\rho>0\), direct one-variable maximization gives \[
\sup_{\psi\in\mathbb R}\{u\psi-\alpha\rho(e^{\psi/\alpha}-1)\}
 =\alpha\{u\log(u/\rho)-u+\rho\}
 =\alpha\rho h(u/\rho).\] The optimizer is \(\psi=\alpha\log(u/\rho)\). Source (3.20), p. 16, and Lemma 4.6, pp. 21–22, contain the \(\rho\)-weighted cumulant in [eq:legendre] (\#eq:legendre). The change-of-measure computation (10.24), p. 45, likewise explicitly contains \(\alpha\rho(x)h(u_0(x)/\rho(x))\), but identifies its integral with \(\alpha\mathcal H_\rho(u_0)\). That identification does not follow from [eq:printedH] (\#eq:printedH); the same false identification appears in (B.32), p. 58.\par
\subsection*{A smooth positive failure of Theorem 4}
For a separate direct check, keep \(d=1,\alpha=2,t_{\mathrm{fin}}=1\) and set \[u(t,x)=2+\tfrac t2\cos(2\pi x),\qquad
\theta(t,x)=-\frac{\sin(2\pi x)}{2\pi u(t,x)},\qquad
g=u_x+\tfrac12\theta.\] The path is smooth, has mass \(2\), and satisfies \(3/2\le u\le5/2\). Direct differentiation verifies both equations [eq:skeleton] (\#eq:skeleton) and [eq:action] (\#eq:action), so all regularity and finite-cost hypotheses hold. These controls attain the respective infima, independently of the source’s dual-action formula: any other one-dimensional periodic control differs by \(c(t)/u\), and \[\int_{\mathbb T}\theta/u\,dx=0,\qquad
\int_{\mathbb T}g/u\,dx=0.\] The first equality follows from oddness in \(x\); the second also uses \(\int u_x/u=0\). The cross terms in the squared norms therefore vanish, proving minimality.\par
Equation reference: eq:smoothEDI. Put \(\Delta\mathcal H=\mathcal H(u_1)-\mathcal H(u_0)\) and \(D=\int_0^1\mathcal D_2(u_t)\,dt\). Periodic integration by parts in the continuity equation gives \(2\Delta\mathcal H=\int_0^1\int_{\mathbb T}u_x\theta\). Expanding the norm of \(g=u_x+\theta/2\) yields \[
\mathcal J(u)=\Delta\mathcal H+\tfrac12D+\tfrac14\mathcal A(u).\] Strict convexity of \(h\), and the fact that \(u_1\) is nonconstant of mean \(2\), imply \(\Delta\mathcal H>0\). For constant \(\rho\), mass conservation and [eq:printedH] (\#eq:printedH) give \(\mathcal H_\rho(u_1)-\mathcal H_\rho(u_0)=\Delta\mathcal H/\rho\). At the allowed choice \(\rho=2\), the right-hand side of source (1.19) is consequently smaller than [eq:smoothEDI] (\#eq:smoothEDI) by \(\Delta\mathcal H/2>0\). This refutes the literal identity on an everywhere positive smooth path.\par
\subsection*{Correction, counterevidence and impact}
Equation reference: eq:correctH. The compatible entropy is \[
\widehat{\mathcal H}_\rho(u)=\int_{\mathbb T^d}\rho\,h(u/\rho)\,dx
 =\int_{\mathbb T^d}\{u\log(u/\rho)-u+\rho\}\,dx.\] It restores the pointwise dual [eq:legendre] (\#eq:legendre), both exact event costs, and the smooth calculation above. For constant \(\rho\) on a fixed-mass path, \(\widehat{\mathcal H}_\rho(u_1)-\widehat{\mathcal H}_\rho(u_0)=\mathcal H(u_1)-\mathcal H(u_0)\), as the gradient-flow argument requires.\par
The case \(\rho=1\) has no discrepancy. Interpreting the measure in source (3.6) as \(\rho\,dx\) would also repair the definition, but the actual page displays \(dx\). This is strong evidence of an intended local correction, especially because (10.24) already contains the required weight; it does not validate the formula as printed. Neither an alternative definition of the equilibrium parameter nor a rescaling of the LDP speed is justified by the explicit laws and hypotheses used above. The positive witness also rules out a repair consisting only of excluding the vacuum.\par
The assessment is Major: the affected rate is the central conclusion of Theorem 2 and an input to the time-reversal argument for Theorem 4. The corrected formula is locally verified here; the entire revised LDP and irregular-path identity are not independently reproved. The hydrodynamic equation itself is not refuted.\par
\section*{Problem 2: control and action inconsistencies in the core proof}
Problem 2 (Major for the control derivation; local action discrepancy).  Source (11.9), p. 48, proposes an infeasible optimizing control in the proof of Theorem 4. Related action coefficients in (11.6), (11.17) and (B.33)–(B.34) do not agree with [eq:action] (\#eq:action).\par
\subsection*{The sign of the proposed control}
The beginning of source §11, p. 47, correctly maps a continuity control to a skeleton control by \(g=\nabla u^{\alpha/2}+\theta/2\). Equation (11.9) instead uses \[g_-=\theta/2-\Pi[u]\nabla u^{\alpha/2},\] where \(\Pi[u]\) projects onto the closure of \(u^{\alpha/2}\nabla\varphi\). This is asserted to be feasible and optimal.\par
Take \(\alpha=2\) and \(u(t,x)=v(x)=2+\cos(2\pi x)\) on \([0,1]\times\mathbb T\). The path is smooth, strictly positive and stationary. The continuity minimizer is \(\theta=0\), since it is feasible with zero cost. Also \(v_x=v(\log v)_x\), so the projection fixes \(v_x\). The printed choice is therefore \(g_-=-v_x\), for which the skeleton right-hand side equals \[\tfrac12(v^2)_{xx}-(vg_-)_x=(v^2)_{xx}.\] At \(x=0\) this is \(-24\pi^2\), whereas \(u_t=0\). The candidate is not feasible. The correct choice \(g_+=v_x\) makes the right-hand side identically zero. This verifies every hypothesis without relying on an irregular PDE existence assertion.\par
\subsection*{An exact action test}
Equation reference: eq:trig. A different positive path detects the coefficient discrepancies: \[u(t,x)=2+\cos(2\pi(x-t)),\qquad
\theta=2-\frac3u,\qquad 0\le t\le1.\] Here \(u_t=-u_x\) and \((u\theta/2)_x=u_x\), so the continuity equation holds. Elementary integration gives \[
\int_{\mathbb T}\frac{dx}{u}=\frac1{\sqrt3},\qquad
\int_{\mathbb T}\frac{dx}{u^2}=\frac2{3\sqrt3},\qquad
\int_{\mathbb T}\frac{\theta}{u}\,dx=0.\] For example, the substitution \(z=\tan(s/2)\) in \((2\pi)^{-1}\int_0^{2\pi}(a+\cos s)^{-1}ds\) gives \((a^2-1)^{-1/2}\); differentiation in \(a>1\) gives the second integral at \(a=2\). Thus these evaluations need no numerical quadrature.\par
Equation reference: eq:actionvalue. As above, other feasible controls differ by \(c(t)/u\), and [eq:trig] (\#eq:trig) proves that \(\theta\) is the unique minimizer. Hence \[
\mathcal A(u)=\tfrac12\int_0^1\int_{\mathbb T}\theta^2=2-\sqrt3>0.\] The correct skeleton minimizer is \(g=u_x+\theta/2\). Since \(\int u_x\theta=0\) and \(\int u_x^2=2\pi^2\), \[\mathcal J(u)=\pi^2+\tfrac14(2-\sqrt3).\] Entropy is constant along this translating path. The literal (11.6) instead gives \(\pi^2+\tfrac12(2-\sqrt3)\), so it is not an identity. The correct coefficient is already visible in source (11.15), p. 49, and in the statement (1.19); it must be retained in (11.6) and (11.17).\par
For the separate dual formula, let \[B_u(\varphi)=\langle\varphi_1,u_1\rangle-\langle\varphi_0,u_0\rangle
 -\int_0^1\langle\partial_t\varphi,u\rangle\,dt .\] The continuity equation gives \(B_u(\varphi)=\tfrac12\langle\theta,u\nabla\varphi\rangle\). Source (B.33)–(B.34), pp. 58–59, claims \(\mathcal A(u)=\sup_\varphi\{B_u(\varphi)-\tfrac12\|u\nabla\varphi\|_2^2\}\). Completing the square shows that this supremum is actually \(\|\theta\|_2^2/8=\mathcal A(u)/4\). It is attained for the present example: \(\varphi_x=\theta/(2u)\) is smooth and has zero spatial mean by [eq:trig] (\#eq:trig), so a smooth periodic \(\varphi\) exists. Thus the displayed dual is \((2-\sqrt3)/4\), not \(2-\sqrt3\).\par
\subsection*{Verified algebraic correction and boundaries}
Write \(w=\Pi[u]\nabla u^{\alpha/2}\). The feasible projected control is \(g=\theta/2+w\), agreeing with the earlier source p. 47. Under time reversal, its companion is \(g^r=Tw-T\theta/2\), so \[g^r+Tg=2Tw,\qquad g^r-Tg=-T\theta.\] With a valid time-reversal rate identity, the parallelogram calculation therefore yields \[\mathcal J(u)=\tfrac12\{\alpha\Delta\widehat{\mathcal H}_\rho+\|w\|_2^2+\mathcal A(u)/2\}.\] The dual of the action [eq:action] (\#eq:action) is normalized by \[\sup_\varphi\{B_u(\varphi)-\tfrac18
 \|u^{\alpha/2}\nabla\varphi\|_2^2\},\] as follows by completing the square with the continuity constraint. The replacement restores the exact smooth examples.\par
The control error is Major at the derivation level: feasibility and optimality feed (11.11)–(11.15), the central passage from reversibility to the claimed entropy–action identity. The correct map on p. 47 is explicit counterevidence against interpreting this as a new obstruction to the intended theorem. The factor-four appendix dual discrepancy is local: multiplying by the correct positive constant preserves the lower-semicontinuity argument’s qualitative conclusion. Neither example disproves the correctly normalized identity at \(\rho=1\), or the existence of minimizing controls. A complete proof for all irregular finite-cost paths is not supplied by these local corrections.\par
\section*{Problem 3: the sharp discrete dissipation bound loses a factor two}
Problem 3 (Minor: quantitative \(\Gamma\)-lower bound).  Source (5.4) and the conclusion \(\mathcal F_\alpha(u)\le z\) in Lemma 5.1, p. 23, are incompatible with the discrete normalization (5.2), p. 22, and the continuum normalization [eq:continuousD] (\#eq:continuousD).\par
Equation reference: eq:lattice. In dimension one, \(\alpha=2\), nearest-neighbour probabilities are \(p^N(k,k\pm1)=1/2\), as also used in source (B.2), p. 53. Thus \[\mathcal D_{2,N}(\eta)=\frac N2\sum_{k=0}^{N-1}
 (\eta_{k+1}-\eta_k)^2.\] Take \(\chi_N=N^{-2}\), and define a time-independent lattice path by \[\eta_k^N=2+\frac{k(N-k)}{N^2},\quad 0\le k<N.\] These are allowed particle values: the particle counts are the nonnegative integers \(2N^2+k(N-k)\). The paths belong to \(\mathbb D_N\) and converge uniformly, hence in \(\mathbb D\), to the time-independent periodic function \[f(x)=2+x(1-x),\quad 0\le x\le1.\] This function is positive and periodic Lipschitz; the derivative jump at the period boundary is allowed by the lemma and \(f\in H^1(\mathbb T)\). Its entropy is finite. Direct finite summation gives \[
\mathcal D_{2,N}(\eta^N)
 =\frac1{2N^3}\sum_{k=0}^{N-1}(N-2k-1)^2
 =\frac{N^2-1}{6N^2}\longrightarrow\frac16,
\qquad
\mathcal D_2(f)=\int_0^1(1-2x)^2dx=\frac13.\] Uniform convergence within a fixed positive compact range also gives \(\mathcal H(\eta^N)\to\mathcal H(f)\). With time horizon one, \[\mathcal F_{2,N}(\eta^N)\to\mathcal H(f)+\tfrac16
 <\mathcal H(f)+\tfrac13=\mathcal F_2(f).\] This is a counterexample to (5.4) with all lattice and energy hypotheses checked. Equivalently, choose \(z=\mathcal H(f)+1/4\); the required bound holds for all sufficiently large \(N\), and finitely many earlier terms may be replaced by the constant configuration \(2\) without changing the limit.\par
The correct factor is already suggested by the sentence preceding source (5.42), p. 27, which states the limit of \(\mathcal D_{2,N}\) as \(\tfrac12\|\nabla\cdot\|_2^2\). More generally, putting \(\xi_N=(\eta^N)^{\alpha/2}\), (5.2) equals \[\frac1{\alpha N^d}\sum_{x,i}
 |N(\xi_N(x+e_i/N)-\xi_N(x))|^2.\] Strong convergence of \(\xi_N\) in \(L^2\), together with bounded difference quotients, identifies their weak limits as \(\partial_i\xi\). Weak lower semicontinuity then gives \(\tfrac12\mathcal D_\alpha(u)\le\liminf_N\mathcal D_{\alpha,N}(\eta^N)\). The corresponding corrected energy lower bound is \[\sup_t\mathcal H(u_t)+\tfrac12\int_0^1\mathcal D_\alpha(u_t)\,dt
 \le\liminf_N\mathcal F_{\alpha,N}(\eta^N),\] and it implies the weaker \(\mathcal F_\alpha(u)\le2z\) under the original uniform bound \(z\).\par
This is a local quantitative defect. Finiteness of the energy and the qualitative weak-to-strong and exclusion arguments are not refuted: a fixed factor does not change whether the dissipation is finite. Changing the lattice transition convention would also alter the generator’s displayed \(\tfrac12\Delta\) limit and is not an available interpretation of the original formulas.\par
\section*{Problem 4: the critical entropy moment can be finite}
Problem 4 (Minor: endpoint \(\gamma=\alpha\)).  Remark 4.5 and (4.23), p. 20, assert infinite entropy exponential moments as soon as \(\gamma\ge\alpha\). The endpoint assertion is false for allowed constant reference densities \(0<\rho<1\).\par
Equation reference: eq:critical. For a single site with size \(\chi>0\), put \(h(s)=s\log s-s+1\) as above. At \(\gamma=\alpha\), the exact law on p. 13 gives \[
\mathbb E_{\pi_\rho^N}
 \exp\{\alpha h(\chi m)/\chi\}
=\frac{e^{\alpha/\chi}}{Z_\alpha(\rho/\chi)}
 \sum_{m=0}^{\infty}\rho^{\alpha m}
 \left(\frac{m^m e^{-m}}{m!}\right)^\alpha ,\] with the \(m=0\) term equal to one. The elementary estimate \(\log(m!)\ge\int_1^m\log x\,dx\ge m\log m-m\) implies \[\sum_{m=0}^{\infty}\rho^{\alpha m}
 \left(\frac{m^m e^{-m}}{m!}\right)^\alpha
 \le\frac1{1-\rho^\alpha}<\infty .\] The full moment in (4.23) is the product of \(N^d\) finite single-site moments, by independence and \(\mathcal H(\eta)=N^{-d}\sum_xh(\eta(x))\). For instance, \(\alpha=2,\rho=1/2,\chi_N=N^{-2}\) satisfies all standing scaling and positivity conditions and contradicts the asserted divergence for every finite \(N\).\par
For fixed positive constant \(\rho\), Stirling’s formula sharpens the endpoint summand to a constant times \(\rho^{\alpha m}m^{-\alpha/2}\). Thus the critical moment is finite for \(\rho<1\), infinite for \(\rho>1\), and at \(\rho=1\) finite exactly when \(\alpha>2\). For \(\gamma>\alpha\) and \(\rho>0\), the positive \((\gamma-\alpha)m\log m\) term does imply divergence. The case \(\rho=0\) is the deterministic vacuum and must be excluded from any universal divergence statement.\par
Replacing \(h(s)\) by \(s\log s\) changes the critical threshold: the linear term cannot be discarded for this endpoint question. The proof of Lemma 4.4 uses that replacement only for an upper bound at \(\gamma<\alpha\), where the present example does not contradict it. The defect is therefore local to the sharpness remark; the theorem’s strictly subcritical hypothesis is not invalidated.\par
\section*{Problem 5: the mass-uniform prefactor is too small}
Problem 5 (Minor: statement of Lemma 5.3).  Source (5.11), p. 24, claims, for some \(\beta>\alpha>1\) and \(C=C(d,\alpha)\), \[\|\eta^N\|_\beta^\beta
 \le C(1+\langle1,\eta^N\rangle)(1+\mathcal D_{\alpha,N}(\eta^N))\] for all nonnegative lattice functions. No such mass-independent constant exists with this linear mass factor.\par
For any fixed lattice, take \(\eta^N\equiv K\ge1\). It belongs to every required Lebesgue space, its mass is \(K\), and its discrete dissipation is zero. The claim becomes \(K^\beta\le C(1+K)\). Choosing \(K^{\beta-1}>2C\) contradicts it. This works for every \(\beta>\alpha>1\); no choice of the existential exponent repairs the display.\par
The proof itself ends at (5.16) with the different factor \(1+\langle1,\eta^N\rangle^\beta\). This is consistent with the interpolation in (5.14)–(5.15): with \(a=\|\eta^N\|_1\) and \(\beta=\alpha+1-2/p\), those lines give \[\|\eta^N\|_\beta^\beta
 \le C a^{1-2/p}(a^\alpha+\mathcal D_{\alpha,N})
 \le C(1+a^\beta)(1+\mathcal D_{\alpha,N}).\] The last inequality also holds at \(a=0\). This verifies the mass scaling correction and identifies it in the existing proof; it is not a new proof of the external discrete Sobolev input to (5.14). Subsequent fixed-mass uses may absorb the corrected power into a constant depending on the mass cap. Accordingly no qualitative compactness or hydrodynamic failure is established by this local counterexample.\par
\section*{Impact and conclusion}
The relative-entropy error gives direct counterexamples to core statements with their literal definitions. The action/control problem invalidates specific essential equalities in the gradient-flow derivation, while the correct control map and final coefficients are partly present elsewhere in the same source. The remaining defects are local; their examples do not disprove the qualitative weak-to-strong principle, the strictly subcritical moment bound, or estimates on a fixed mass class.\par
Object | Assessment | Verified correction and limit\par
Relative entropy | Major | The weight \(\rho\) restores the dual, exact event costs and smooth entropy identity; the full revised stochastic/PDE theory is not certified.\par
Control/action proof | Major derivation; local dual subfinding | The plus control, coefficient \(\mathcal A/2\) inside the EDI bracket, and dual penalty \(1/8\) restore the exact algebra and smooth tests.\par
Sharp energy bound | Minor | A half-weighted continuum dissipation matches the lattice; no qualitative theorem is refuted by the constant mismatch.\par
Critical moment | Minor | Endpoint convergence depends on \(\rho,\alpha\); the strict subcritical result survives this check.\par
Uniform mass factor | Minor | The proof’s \(1+a^\beta\) factor has the required scaling; external Sobolev inputs were not independently reproved.\par
For every item, author adoption is unknown. Scientific adjudication is pending and no human review or independent acceptance is asserted. The local repairs are evidence about specified formulas and examples, not certification of the whole paper after revision.\par
\section*{Supporting checks}
cd4216b8c8ff425778dbae0ed00c2edf4691914fd7bb2fc75e8fdaa66252837b.\par
Actual page-image inspection confirmed the displayed formulas, their hypotheses and competing normalizations on source pp. 1–6, 12–13, 16–17, 20–24, 27, 45, 47–50, 53 and 58–59. The detailed coverage ledger distinguishes these image views from the complete text reading.\par
The independently written `support/exact\_checks.py` records the actual interpreter and package versions, source/script hashes and outputs. Exact SymPy checks verify the scalar Legendre identity, the two control residuals, the smooth entropy-path constraints, the translating-path constraint and orthogonality, and the finite lattice sum. The limiting arguments and minimizer proofs are analytic arguments in this report, not theorems certified by software. The supplemental 80-digit calculations are finite diagnostics: \[\begin{array}{c|cc}
N & -N^{-2}\log Z_2(N^2/2) &
 N^{-2}\log\pi_{1/2}^N(1)\\
\hline
4&-0.8564275894&-0.5315270193\\
64&-0.9987603049&-0.3875340810\\
\text{analytic limit}&-1&1-2\log2
\end{array}\] The second column concerns the zero-path event after extracting the full spatial speed; the third concerns the positive initial-profile event. Neither finite table is used in place of the limit proof. The same output records the exact action \(2-\sqrt3\), the discrete dissipation limit \(1/6\) versus \(1/3\), and a positive numerical entropy-identity gap \(0.01574969066\ldots\), whose strict sign was proved analytically.\par
The existing source extraction was reused with its original attribution. No author code was executed. Missing optional code was not treated as a mathematical defect. Actual compilation, parsing, rendering and subsequent page views are bound to the current report in the matching verification records. No external Chat, separate model selection, backend equivalence, author response or institutional endorsement is claimed.\par
9 B. Gess and D. Heydecker, The porous medium equation: Large deviations and gradient flow with degenerate and unbounded diffusion, arXiv:2303.11289v3, 24 March 2025, 62 pp., https://arxiv.org/abs/2303.11289v3. Associated publication: Comm. Pure Appl. Math. 78 (2025), no. 9, 1609–1655, https://doi.org/10.1002/cpa.22251. The mathematical evidence in this report is the acquired preprint.\par
\end{document}
